% SOURCE_STATEMENT: "For $A(X_r)=-\\beta^{-1}\\log\\int_{\\mathbb R^{dN_c}}e^{-\\beta V(X_r,X_c)}\\,dX_c$, can $A(X_r)/N_c$ be computed efficiently as $N_c\\to+\\infty$?"
% Preparation and final self-audit: CONTRIBUTING.md prompts C and D; source/audit artifacts retained in the local physics_statements_2026_10_11 batch.
% !TeX program = lualatex
\documentclass[11pt]{article}
\usepackage[margin=25mm]{geometry}
\usepackage{fontspec}
\setmainfont{Latin Modern Roman}
\usepackage{amsmath,amssymb,mathtools}
\usepackage{unicode-math}
\setmathfont{Latin Modern Math}
\usepackage{xurl,hyperref}
\hypersetup{colorlinks=true,urlcolor=blue}
\setcounter{secnumdepth}{0}
\setlength{\parindent}{0pt}
\setlength{\parskip}{5pt}
\setlength{\emergencystretch}{3em}
\newcommand{\sref}[2]{\hyperlink{source:#1:#2}{[S#2]}}
\newcommand{\cataloguescope}[1]{\paragraph{Recorded target.}#1}
\newcommand{\R}{\mathbb R}
\newcommand{\C}{\mathbb C}
\newcommand{\N}{\mathbb N}
\newcommand{\Z}{\mathbb Z}
\newcommand{\Q}{\mathbb Q}
\newcommand{\F}{\mathbb F}
\newcommand{\eps}{\varepsilon}
\begin{document}
\section{Efficient Reduced Free-Energy Computation for Atomistic Systems}\label{statement}

\subsection{Definitions and mathematical statement}

\paragraph{Shared notation.}$\mathbb N=\{1,2,\ldots\}$; the standard number fields have their usual meanings. All additional conventions are stated locally.


\paragraph{Definitions and conventions.} $X_r\in\mathbb R^{dN_r}$ are retained representative-atom coordinates and $X_c\in\mathbb R^{dN_c}$ are the coordinates to be integrated out; $N=N_r+N_c$ and $\beta>0$ is inverse temperature. $V$ is the atomistic potential energy and the integral is conditional on the fixed retained coordinates. The normalized free energy $A(X_r)/N_c$ is measured per eliminated particle. Existence of a finite integral requires confinement or the boundary and translation-fixing conventions of the atomistic model.

\paragraph{Source scope and qualifications.} Efficiency and the allowed potentials are not quantitatively specified. The report describes a method for chain-like nearest- and next-nearest-neighbor systems even when atom coordinates are multidimensional. Its passage closely repeats ID 30001115; the latter clean formulation explicitly restricts to a one-dimensional system, so the potentially broader catalogue record is retained with this relation disclosed.

\paragraph{Catalogue mathematical target.}
\paragraph{Catalogue wording retained for traceability.}
For $A(X_r)=-\beta^{-1}\log\int_{\mathbb R^{dN_c}}e^{-\beta V(X_r,X_c)}\,dX_c$, can $A(X_r)/N_c$ be computed efficiently as $N_c\to+\infty$?

\paragraph{Saved literature context.} The supplied catalogue's assessment dated 2026-08-21 records: The extracted question asks for an efficient large-Nc computation and no general algorithm or asymptotic resolution is identified in the cited literature. The assessment is retained as dated provenance; it is not a new comprehensive literature or open-status certification.
\subsection{Short English statement}
Compute the reduced free energy per removed atom efficiently as the number of removed atoms grows.
\subsection{Sources}
\begin{itemize}
\item[\hypertarget{source:statement:1}{S1}] ULAM.AI and the UnsolvedMath Contributors, UnsolvedMath, catalogue entry. Catalogue curation under CC BY 4.0; definitions and English exposition added locally.
\url{https://huggingface.co/datasets/ulamai/UnsolvedMath}.
\item[\hypertarget{source:statement:2}{S2}] Computational Multiscale Methods (2009), Oberwolfach Reports, DOI 10.4171/owr/2009/29, pp. 1597--1666; official report source and mathematical context.
\url{https://doi.org/10.4171/owr/2009/29}.
\item[\hypertarget{source:statement:3}{S3}] Official EMS Press PDF of the same Oberwolfach report; PDF page 26.
\url{https://ems.press/content/serial-article-files/46232}.
\end{itemize}
\label{end:statement}
\end{document}
